\documentclass[12pt]{article}

\usepackage{amsmath, amssymb, amsthm}
\usepackage{geometry}
\usepackage{graphicx}
\usepackage{setspace}
\usepackage{titlesec}

\geometry{margin=1in}
\setstretch{1.2}

\title{\textbf{Lecture 18: Marginal PMF, Correlation, Covariance, Joint Gaussian Random Variables}}
\author{Lecture Scribe}
\date{}

\begin{document}

\maketitle

\section*{1. Marginal PMF}

\subsection*{1.1 Definition}

Consider two discrete random variables $X$ and $Y$.

The joint PMF is defined as

\[
p_{XY}(x,y) = P(X=x, Y=y)
\]

This gives the probability that both events occur simultaneously.

However, often we want the probability distribution of only one variable.

To obtain it, we sum the joint probabilities over all possible values of the other variable.

Thus,

\[
p_X(x) = \sum_{y} p_{XY}(x,y)
\]

\[
p_Y(y) = \sum_{x} p_{XY}(x,y)
\]

\textbf{Key Idea}

The marginal PMF of one variable is obtained by summing the joint PMF over all possible values of the other variable.

\subsection*{1.2 Marginal PMF from a Joint PMF Table}

Suppose the joint PMF is given as

\[
\begin{array}{c|ccc|c}
X \backslash Y & y_1 & y_2 & y_3 & p_X(x) \\
\hline
x_1 & p_{11} & p_{12} & p_{13} & \sum_j p_{1j} \\
x_2 & p_{21} & p_{22} & p_{23} & \sum_j p_{2j} \\
x_3 & p_{31} & p_{32} & p_{33} & \sum_j p_{3j} \\
\hline
p_Y(y) & \sum_i p_{i1} & \sum_i p_{i2} & \sum_i p_{i3} &
\end{array}
\]

Thus

\[
p_X(x_i)=\sum_j p_{ij}
\]

\[
p_Y(y_j)=\sum_i p_{ij}
\]

Important property

\[
\sum_i p_X(x_i)=1
\]

\[
\sum_j p_Y(y_j)=1
\]

Thus the marginals are valid probability distributions.

\subsection*{1.3 Example: Computing Marginal PMFs}

Suppose the joint PMF is

\[
\begin{array}{c|cc}
 & y=0 & y=1 \\
\hline
x=0 & \frac{1}{8} & \frac{1}{4} \\
x=1 & \frac{3}{8} & \frac{1}{4}
\end{array}
\]

\textbf{Step 1: Compute $p_X(x)$}

For $x=0$

\[
p_X(0)=\frac{1}{8}+\frac{1}{4}=\frac{1}{8}+\frac{2}{8}=\frac{3}{8}
\]

For $x=1$

\[
p_X(1)=\frac{3}{8}+\frac{1}{4}=\frac{3}{8}+\frac{2}{8}=\frac{5}{8}
\]

Thus

\[
p_X(x)=
\begin{cases}
\frac{3}{8}, & x=0 \\
\frac{5}{8}, & x=1
\end{cases}
\]

\textbf{Step 2: Compute $p_Y(y)$}

For $y=0$

\[
p_Y(0)=\frac{1}{8}+\frac{3}{8}=\frac{4}{8}=\frac{1}{2}
\]

For $y=1$

\[
p_Y(1)=\frac{1}{4}+\frac{1}{4}=\frac{1}{2}
\]

Thus

\[
p_Y(y)=
\begin{cases}
\frac{1}{2}, & y=0 \\
\frac{1}{2}, & y=1
\end{cases}
\]

\section*{2. Correlation}

\subsection*{2.1 Definition}

A basic measure of the joint behavior of two variables is the product moment

\[
R_{XY}=E[XY]
\]

For continuous random variables

\[
R_{XY}=\int_{-\infty}^{\infty}\int_{-\infty}^{\infty} xy f_{XY}(x,y) \, dx\, dy
\]

This represents the average value of the product $XY$.

It is often called a raw correlation measure.

\subsection*{2.2 Limitations}

The raw product moment has some issues:

\begin{itemize}
\item It depends on the scale of $X$ and $Y$.
\item It is computed about the origin, not about the mean.
\end{itemize}

Therefore it mixes together

\begin{itemize}
\item mean values
\item spread
\item dependence
\end{itemize}

Thus it is not a clean measure of centered dependence.

To remove the effect of means, we center the variables.

This leads to covariance.

\section*{3. Covariance}

\subsection*{3.1 Definition}

Covariance measures how two variables vary together around their means.

\[
\text{Cov}(X,Y)=E[(X-\mu_X)(Y-\mu_Y)]
\]

where

\[
\mu_X=E[X], \quad \mu_Y=E[Y]
\]

\textbf{Interpretation}

\begin{enumerate}
\item Subtract the mean from each variable
\item Measure how their deviations vary together
\end{enumerate}

\subsection*{3.2 Interpretation}

\textbf{Positive Covariance}

\[
\text{Cov}(X,Y) > 0
\]

When one variable increases, the other tends to increase.

Example: GPA vs hours studied.

\textbf{Negative Covariance}

\[
\text{Cov}(X,Y) < 0
\]

When one increases, the other tends to decrease.

Example: speed vs travel time.

\textbf{Zero Covariance}

\[
\text{Cov}(X,Y)=0
\]

No linear co-variation trend.

Important note:

Zero covariance means uncorrelated but does not necessarily imply independence.

\section*{4. Algebraic Form of Covariance}

Start with the definition

\[
\text{Cov}(X,Y)=E[(X-\mu_X)(Y-\mu_Y)]
\]

Expand

\[
(X-\mu_X)(Y-\mu_Y)=XY-X\mu_Y-Y\mu_X+\mu_X\mu_Y
\]

Taking expectation

\[
E[XY-X\mu_Y-Y\mu_X+\mu_X\mu_Y]
\]

Using linearity of expectation

\[
E[XY]-\mu_YE[X]-\mu_XE[Y]+\mu_X\mu_Y
\]

Since

\[
E[X]=\mu_X, \quad E[Y]=\mu_Y
\]

Substitute

\[
E[XY]-\mu_X\mu_Y-\mu_X\mu_Y+\mu_X\mu_Y
\]

Thus

\[
\boxed{\text{Cov}(X,Y)=E[XY]-E[X]E[Y]}
\]

\section*{5. Properties of Covariance}

\begin{enumerate}
\item Symmetry
\[
\text{Cov}(X,Y)=\text{Cov}(Y,X)
\]

\item Variance as special case
\[
\text{Cov}(X,X)=\text{Var}(X)
\]

\item Linearity with constants
\[
\text{Cov}(aX+b,cY+d)=ac\,\text{Cov}(X,Y)
\]

\item Independence

If $X \perp Y$

\[
\text{Cov}(X,Y)=0
\]

But the reverse is not always true.
\end{enumerate}

\section*{6. Example 1}

Given

\[
Y=-2X+3
\]

Find $\text{Cov}(X,Y)$.

Step 1

\[
\text{Cov}(X,Y)=E[XY]-E[X]E[Y]
\]

Step 2

\[
XY=X(-2X+3)
\]

\[
XY=-2X^2+3X
\]

\[
E[XY]=E[-2X^2+3X]
\]

Step 3

\[
E[Y]=-2E[X]+3
\]

Step 4 Substitute

\[
\text{Cov}(X,Y)=E[-2X^2+3X]-E[X](-2E[X]+3)
\]

Simplifying

\[
\text{Cov}(X,Y)=-2E[X^2]+2(E[X])^2
\]

But

\[
\text{Var}(X)=E[X^2]-(E[X])^2
\]

Thus

\[
\boxed{\text{Cov}(X,Y)=-2\text{Var}(X)}
\]

\section*{7. Pearson Correlation Coefficient}

Covariance depends on units.

To normalize we divide by standard deviations.

\[
\rho_{XY}=\frac{\text{Cov}(X,Y)}{\sqrt{\text{Var}(X)\text{Var}(Y)}}
\]

Properties

\begin{itemize}
\item Measures strength and direction of linear relationship
\item Unit-free
\item Normalized covariance
\end{itemize}

Range

\[
-1 \leq \rho_{XY} \leq 1
\]

Interpretation

\begin{itemize}
\item $\rho>0$ positive linear trend
\item $\rho<0$ negative linear trend
\item $\rho=0$ no linear correlation
\end{itemize}

\section*{8. Example 2 (Joint PDF)}

Given joint density

\[
f_{XY}(x,y)=
\begin{cases}
\frac{xy}{96}, & 0<x<4, \; 1<y<5 \\
0, & \text{otherwise}
\end{cases}
\]

Compute

\begin{enumerate}
\item $E[X]$
\item $E[Y]$
\item $E[XY]$
\item $E[2X+3Y]$
\item $\text{Var}(X)$
\item $\text{Var}(Y)$
\item $\text{Cov}(X,Y)$
\item $\rho_{XY}$
\end{enumerate}

Step 1

\[
E[X]=\int_0^4 \int_1^5 x f_{XY}(x,y)\, dy\, dx
\]

Substituting

\[
E[X]=\int_0^4 \int_1^5 x\frac{xy}{96} \, dy\, dx
\]

Result

\[
E[X]=\frac{3}{8}
\]

Step 2

\[
E[Y]=\frac{9}{31}
\]

Step 3

Since $X$ and $Y$ are independent

\[
E[XY]=E[X]E[Y]
\]

\[
E[XY]=\frac{3}{8}\times\frac{9}{31}=\frac{27}{248}
\]

Step 4

\[
E[2X+3Y]=2E[X]+3E[Y]
\]

Step 5

\[
\text{Var}(X)=\frac{9}{8}
\]

Step 6

\[
\text{Var}(Y)=\frac{81}{92}
\]

Step 7

\[
\text{Cov}(X,Y)=E[XY]-E[X]E[Y]
\]

\[
\text{Cov}(X,Y)=0
\]

Step 8

\[
\rho_{XY}=\frac{\text{Cov}(X,Y)}{\sqrt{\text{Var}(X)\text{Var}(Y)}}=0
\]

\section*{9. Jointly Gaussian Random Variables}

\subsection*{Definition}

A pair $(X,Y)$ is jointly Gaussian if

\begin{itemize}
\item Their joint density is bivariate Gaussian
\item Each marginal is Gaussian
\end{itemize}

\[
X \sim N(\mu_X,\sigma_X^2)
\]

\[
Y \sim N(\mu_Y,\sigma_Y^2)
\]

\subsection*{Joint Gaussian PDF}

\[
f_{XY}(x,y)=
\frac{1}{2\pi\sigma_X\sigma_Y\sqrt{1-\rho^2}}
\exp
\left(
-\frac{1}{2(1-\rho^2)}
\left[
\left(\frac{x-\mu_X}{\sigma_X}\right)^2
+
\left(\frac{y-\mu_Y}{\sigma_Y}\right)^2
-
2\rho
\left(\frac{x-\mu_X}{\sigma_X}\right)
\left(\frac{y-\mu_Y}{\sigma_Y}\right)
\right]
\right)
\]

where

\[
\rho = \text{correlation coefficient}
\]

\subsection*{Real Life Example}

Consider climate system variables

\begin{itemize}
\item $X$ : atmospheric CO$_2$ concentration (ppm)
\item $Y$ : global temperature anomaly ($^\circ$C)
\end{itemize}

Environmental randomness may cause their joint behavior to follow a Gaussian model.

If

\[
\rho>0
\]

then higher CO$_2$ levels tend to occur with higher temperatures.

\end{document}